\subsection{Clifford and T with Measurement}

Clifford + T computation are performed in the following three steps.
\begin{enumerate}
	\item Prepare the states $\ket{0}^{\otimes n}$,
	\item Perform a sequence of Clifford + T gates,
	\item Measure in $Z$ basis.
\end{enumerate}

Since we will eventually measure in the $Z$ basis, we have the following observation

\begin{remark}
	Postpending a $\tikzfig{Z-dual}$ to any ancilla quibit will not effect the final measurement result, as we are measuring in the $Z$ basis.
\end{remark}

\subsubsection{Pauli and Clifford}

If we have an initial state of $\ket{T}$, we can perform a $T$ gate on any part of the circuit by a Pauli measurement $\Pi^0 (Z \otimes Z \otimes \vec{I})$ in the following way.

\begin{align}\label{eq:embed-S}
	\tikzfig{q-2-3-1}
\end{align}

Thus for a Clifford + T circuit, we can let the Clifford gates remain the same and replace each $T$ gate with the above Pauli measurement. 
The number of $\ket{T}$ required is exactly equal to the number of $T$ gates in the original circuit.
If there are two consecutive $T$ gates, we can add an identity clifford between them. 
In this way, we have simulated a Clifford + T circuit with alternating Pauli measurements and Clifford gates.

\subsubsection{Only Pauli}

Next, we would like to show that we can simulate any Clifford with Pauli measurement.
We demonstrate this in several steps.

\begin{lemma}[S state with measurement]\label{lm:Pauli-S}
	We can simulate $S$ state by two $Z$ measurements and two $\ket{T}$ in the followng way
	\begin{equation}
		\tikzfig{q-2-3-S-state}
	\end{equation}
\end{lemma}

\begin{lemma}[Multiple of $\frac{\pi}{4}$ state]\label{lm:mul-Z-state}
	Using $k$ $Z$ measurements and $k$ $\ket{T}$, we can simulate the state $\ket{\frac{k\pi}{4}} $ in the exact same way as in lemma \ref{lm:Pauli-S}.
	Importantly, it means we can 
	\begin{enumerate}
		\item simulate a $\tikzfig{Z-state}$ with 8 $Z$ measurements and 8 $\ket{T}$,
		\item simulate a $\tikzfig{Z-Pi-state}$ with 4 $Z$ measurements and 4 $\ket{T}$.
	\end{enumerate}
\end{lemma}


\begin{lemma}[State to Process]
	If we have an ancilla quibit in state $\tikzfig{alpha-state}$, we can embed a $\tikzfig{alpha-process}$ in to the circuit with a $Z\otimes Z \otimes I$ measurement.
\end{lemma}

\begin{proof}
	This can be done in the exact same way as in \eqref{eq:embed-S}, by replacing the $\ket{T}$ state with $\tikzfig{alpha-state}$.	
\end{proof}


\begin{lemma}
	Using 9 $Z$ measurements and 8 ancilla $\ket{T}$, we can simulate the state $\tikzfig{X-state}$ in the ancilla qubit.
\end{lemma}

\begin{proof}
	By	lemma \ref{lm:mul-Z-state}, we can simulate $\tikzfig{Z-state}$ with 8 $Z$ measurements and 8 $\ket{T}$.

	Now use the following $Z$ measurement
	$$
	\tikzfig{q-2-3-sim-X-state}
	$$
\end{proof}
